\section{A Geometric Interpretation of Short- and Long-Term Memory}
\label{sec:interpretation}

We analyze how the geometric structure of $\mathrm{SMN}(n,m)$ determines its long- and short-term memory properties.  The reader may find it helpful to keep the tetrahedron case $\mathrm{SMN}(3,m)$ in mind as a concrete reference throughout.

Long-term memory is encoded along the narrow backbone $\mathcal{B}_{n,m} := \{\alpha \in V_{n,m} : \alpha_2 = \cdots = \alpha_n = 0\}$, a chain of $m$ nodes with potentials $H = 0, 2, \ldots, 2(m{-}1)$ from $a$ to $b$.  Its width is always $1$ while its depth scales with $m$, and skip edges of type $a \to b$ allow this channel to bypass the interior entirely.

Short-term memory resides in $F_{\mathrm{mid}} = F_{\mathrm{in}} \cap F_{\mathrm{out}}$, an isopotential buffer at $H = m-1$ whose cardinality $\binom{m+n-3}{n-2}$ grows with $n$ at fixed depth.  Since $F_{\mathrm{mid}} \subset F_{\mathrm{out}}$, every buffer node is immediately readable at the output interface, making $F_{\mathrm{mid}}$ wide and shallow where the backbone is narrow and deep.

In forward propagation, the backbone and the buffer influence each other's expression through the simplex interior.  Pre-midpoint backbone nodes send activations of type $a \to c_k$ into $F_{\mathrm{mid}}$, so the short-term output is modulated by whatever long-term state the backbone carries.  Symmetrically, $F_{\mathrm{mid}}$ feeds into the post-midpoint backbone via edges of type $c_k \to b$, so the long-term output at $b$ depends on the current short-term buffer.  Each half of the architecture conditions the other's readout at inference time.

In backward propagation, gradients flow in the reverse direction and couple the encoding of the two memory components.  An error signal at the output reaches the backbone directly through the $ab$-skip path, but also reaches $F_{\mathrm{mid}}$ through the post-midpoint interior; the buffer's weights are therefore shaped by the same objective that trains the long-term pathway.  Conversely, gradients propagating back through $F_{\mathrm{mid}}$ reach the pre-midpoint backbone via the $c_k \to b$ reverse path, so long-term encoding is updated in response to how well the short-term buffer performed.  Over training, this mutual gradient flow causes the two components to co-adapt: patterns that recur in the short-term buffer are progressively absorbed into the backbone, while the backbone's evolving state continuously reshapes what the buffer learns to store.
